\section{Introduction}
\label{sec:introduction}

\PARstart{H}{ybrid} vertical take-off and landing (VTOL) aircraft
combine the hovering capability of multirotors with the range and
efficiency of fixed-wing flight, and have consequently become one of
the fastest-growing classes of unmanned aerial vehicles
(UAVs)~\cite{saeed2018survey,ducard2021review}. The price of this
versatility is a control problem that changes structure between flight
regimes, and the literature increasingly treats each regime as a
control problem in its own right~\cite{hu2024hybrid}. This article
focuses on the \emph{fixed-wing flight mode} of a hybrid VTOL platform
--- the mode in which the vehicle spends most of its mission --- and on
the complete path from controller design to embedded validation;
hover and transition are outside its scope.

\subsection{Related Work}
Autopilots for small UAVs remain dominated by nested
proportional--integral--derivative (PID) loops, a
legacy of their simplicity and field-proven
robustness~\cite{chao2010autopilots}, and this architecture underpins
the open-source flight stacks in widespread use
today~\cite{ebeid2018survey,meier2015px4}. Linear--quadratic (LQ)
methods are the natural model-based alternative, and comparisons
between the two families have accumulated since the seminal indoor
quadrotor study of Bouabdallah \emph{et
al.}~\cite{bouabdallah2004pid}: on quadrotor
platforms~\cite{argentim2013pid,rinaldi2023comparative}, on a
commercial mini-drone with experimental PID/LQR/MPC
benchmarks~\cite{okasha2022design}, and, less frequently, on
fixed-wing longitudinal dynamics~\cite{ingabire2019control}.

Hardware-in-the-loop (HIL) simulation --- preceded in practice by
software-in-the-loop (SIL) verification on the same models --- is the
established way to close the gap between simulation and flight: the embedded control system is
exercised against a simulated plant, so faults in the embedded
implementation surface on the bench rather than in the
air~\cite{jung2007hil,cai2009design}. The practice extends from
classical avionics benches to recent microcontroller-class autopilot
platforms~\cite{coopmans2015software,laabousse2025advanced}, is
integral to model-based development workflows in which the flight code
is generated automatically from the verified
model~\cite{erkkinen2006automatic,erkkinen2009model,paw2011integrated},
and has been applied to hybrid VTOL vehicles in the Brazilian research
context~\cite{silva2019control}. Within our research group, this
lineage runs through the Pegasus autopilot HIL
architecture~\cite{bittar2012pegasus}, X-Plane-based HIL of small-UAV
attitude control under atmospheric
disturbances~\cite{bittar2014hil}, and the bench-validated
fixed-wing autopilot of Santos~\cite{santos2018,santos2017hilloops},
whose cascade methodology this work adopts. The aircraft models in
this line of work are identified from flight-test
data~\cite{sato2021,sato2020identificacao,angelo2026}.

\subsection{Research Gaps}
Three gaps recur in this comparative literature:
\begin{enumerate}
  \item \emph{Platform.} Most PID--LQ comparisons target rotary-wing
        platforms; the fixed-wing mode of a hybrid VTOL vehicle is
        rarely the object of study.
  \item \emph{Mission-level embedded evaluation.} The comparisons are
        usually confined to simulation of isolated channels, rather
        than a full mission exercising the longitudinal and
        lateral--directional loops together on the embedded
        implementation.
  \item \emph{Robustness.} Arguably where the two design philosophies
        should differ most, robustness is seldom part of the
        protocol, even though LQ designs carry classical margin
        guarantees~\cite{safonov1977gain} that do not survive output
        feedback in general~\cite{doyle1978guaranteed}.
\end{enumerate}

\subsection{Contributions}
The present work addresses these gaps, extending the group lineage in
three directions: the plant
is the fixed-wing mode of a \emph{hybrid VTOL} aircraft rather than a
conventional fixed-wing one; the embedded controller is not hand-coded
but \emph{automatically generated} from the same Simulink model that
was verified in simulation, so that the tested artifact is the
production code itself; and two different control laws --- a cascade
PID and an output-feedback LQR (LQRy) --- are designed against a
common, literature-grounded requirements set, stressed under model
mismatch and wind, and deployed on the same bench under identical
conditions.

The main contributions of this work are:
\begin{itemize}
  \item the design of a cascade PID architecture for the fixed-wing
        flight mode of a hybrid VTOL aircraft, evaluated against a
        fixed-gain output-feedback LQR (LQRy) designed for the same
        platform, both judged against a common set of
        literature-grounded closed-loop requirements;
  \item a robustness study that inverts the nominal verdict: a
        severe, sign-changing lateral--directional model mismatch
        makes the aggressively tuned cascade PID lose the mission
        outright, while the soft fixed-gain LQRy completes it with
        near-nominal error --- with the remaining scenarios adding
        nuance: under a longitudinal-only mismatch the LQRy stays
        nominal while the PID's clamped altitude loop drifts, under
        wind and gusts the PID is the firmer law, and both tolerate a
        plausible $+10\%$ identification uncertainty;
  \item a low-cost HIL architecture based on
        automatic code generation and a lockstep serial protocol, whose
        SIL--HIL residual is shown to reach the quantization floor of
        the communication bus for the synchronized PID deployment;
  \item the characterization of the lockstep HIL loop as a
        \emph{deterministic} system --- demonstrated by bit-exact
        reproduction of full 300-second flight cases across two
        microcontroller models --- and its use to quantify how much each
        implementation choice affects SIL--HIL agreement; and
  \item a controlled, like-for-like comparison of the two controllers on
        the same rig, same references, and same metrics, with the known
        asymmetries of the two design harnesses disclosed explicitly.
\end{itemize}

\subsection{Organization}
Section~\ref{sec:model} presents the nonlinear model, the trim condition,
and the linearization. Section~\ref{sec:control} details the closed-loop
requirements and both controller designs. Section~\ref{sec:sil} reports
the SIL simulations and the controller comparison.
Section~\ref{sec:robustness} presents the robustness study.
Section~\ref{sec:hil} describes the HIL development, the determinism and
sensitivity analyses, and the embedded results.
Section~\ref{sec:conclusion} concludes the article.
